
\documentclass{article}
\usepackage{colortbl}
\usepackage{makecell}
\usepackage{multirow}
\usepackage{supertabular}

\begin{document}

\newcounter{utterance}

\centering \large Interaction Transcript for game `hot\_air\_balloon', experiment `air\_balloon\_survival\_en\_reasoning off\_hard', episode 1 with llama{-}3.3{-}70b{-}instruct{-}t1.0.
\vspace{24pt}

{ \footnotesize  \setcounter{utterance}{1}
\setlength{\tabcolsep}{0pt}
\begin{supertabular}{c@{$\;$}|p{.15\linewidth}@{}p{.15\linewidth}p{.15\linewidth}p{.15\linewidth}p{.15\linewidth}p{.15\linewidth}}
    \# & $\;$A & \multicolumn{4}{c}{Game Master} & $\;\:$B\\
    \hline

    \theutterance \stepcounter{utterance}  
    & & \multicolumn{4}{p{0.6\linewidth}}{
        \cellcolor[rgb]{0.9,0.9,0.9}{
            \makecell[{{p{\linewidth}}}]{
                \texttt{\tiny{[P1$\langle$GM]}}
                \texttt{Stai partecipando a un gioco di negoziazione collaborativa.} \\
\\ 
\texttt{Insieme a un altro partecipante, dovete accordarvi su un unico set di oggetti da conservare. Ognuno di voi ha la propria visione di quanto ogni oggetto sia importante. Non sai come l'altro partecipante valuta gli oggetti. Inoltre, ti viene fornito lo sforzo che ogni oggetto richiede.} \\
\texttt{Potete accordarvi su un set solo se lo sforzo totale degli oggetti selezionati non supera un limite condiviso:} \\
\\ 
\texttt{LIMITE = 3357} \\
\\ 
\texttt{Ecco i valori di sforzo dei singoli oggetti:} \\
\\ 
\texttt{Sforzo oggetto = \{"A21": 506, "B93": 492, "A11": 338, "B33": 122, "C08": 131, "B39": 909, "B96": 144, "C85": 716, "A82": 263, "B31": 231, "B20": 91, "A60": 651, "B09": 552, "C81": 849, "B21": 720\}} \\
\\ 
\texttt{Ecco la tua visione personale sull'importanza di ogni oggetto:} \\
\\ 
\texttt{Valori di importanza dell'oggetto = \{"A21": 138, "B93": 583, "A11": 868, "B33": 822, "C08": 783, "B39": 65, "B96": 262, "C85": 121, "A82": 508, "B31": 780, "B20": 461, "A60": 484, "B09": 668, "C81": 389, "B21": 808\}} \\
\\ 
\texttt{Obiettivo:} \\
\\ 
\texttt{Il tuo obiettivo è negoziare un set condiviso di oggetti che ti avvantaggi il più possibile (cioè, massimizzi l'importanza totale per TE), rimanendo entro il LIMITEE. Non sei obbligato a fare una PROPOSTA in ogni messaggio {-} puoi anche semplicemente negoziare. Tutte le tattiche sono permesse!} \\
\\ 
\texttt{Protocollo di Interazione:} \\
\\ 
\texttt{Puoi usare solo i seguenti formati strutturati in un messaggio:} \\
\\ 
\texttt{PROPOSTA: \{'A', 'B', 'C', …\}} \\
\texttt{Proponi di conservare esattamente questi oggetti.} \\
\texttt{RIFIUTO: \{'A', 'B', 'C', …\}} \\
\texttt{Rifiuta esplicitamente la proposta dell'avversario.} \\
\texttt{ARGOMENTAZIONE: \{'...'\}} \\
\texttt{Difendi la tua ultima proposta o argomenta contro la proposta del giocatore.} \\
\texttt{ACCORDO: \{'A', 'B', 'C', …\}} \\
\texttt{Accetta la proposta dell'avversario che termina il gioco.} \\
\\ 
\\ 
\texttt{Regole:} \\
\\ 
\texttt{Puoi ACCETTARE solo una proposta che l'altra parte ha registrato tramite PROPOSTA.} \\
\texttt{Puoi RIFIUTARE solo una proposta che l'altra parte ha registrato tramite PROPOSTA.} \\
\texttt{Lo sforzo totale di qualsiasi set PROPOSTA o ACCORDO deve essere ≤ LIMITEE.} \\
\texttt{NON rivelare i tuoi punteggi di importanza nascosti.} \\
\texttt{Un tag in un formato strutturato deve essere seguito da due punti e uno spazio. L'argomento deve essere un set python contenente 0 o più stringhe.} \\
\texttt{Quindi, deve essere della forma TAG: \{...\}} \\
\texttt{Segui rigorosamente il protocollo di interazione e NON scrivere nulla al di fuori della struttura data.} \\
\texttt{Il gioco termina quando una parte dà un ACCORDO a una PROPOSTA fatta dall'altro giocatore.} \\
\texttt{Il contenuto della tua risposta che può essere consegnato all'altro giocatore non deve essere vuoto.} \\
\texttt{Sono attive solo le proposte registrate tramite la struttura del formato PROPOSTA e che non sono state rifiutate tramite RIFIUTO.} \\
\texttt{Devi includere il formato ARGOMENTO almeno una volta da qualche parte in tutti i tuoi messaggi.} \\
            }
        }
    }
    & \\ \\

    \theutterance \stepcounter{utterance}  
    & \multicolumn{4}{p{0.6\linewidth}}{
        \cellcolor[rgb]{0.9,0.9,0.9}{
            \makecell[{{p{\linewidth}}}]{
                \texttt{\tiny{[P1$\rangle$GM]}}
                \texttt{ARGOMENTAZIONE: \{'A21', 'B93', 'A11', 'B33', 'C08', 'B96', 'A82', 'B31', 'B20'\} } \\
\texttt{Sulla base dei valori di sforzo forniti, dobbiamo considerare attentamente quali oggetti includere per massimizzare l'importanza totale rimanendo entro il limite di sforzo. È fondamentale trovare un equilibrio tra oggetti con alti valori di importanza e oggetti con bassi valori di sforzo per ottimizzare la nostra selezione. La proposta iniziale include una combinazione di oggetti con bassi e alti valori di sforzo per coprire una gamma di possibilità senza superare il limite di sforzo. Sarebbe utile conoscere la prospettiva dell'altro partecipante per raffinare la selezione e raggiungere un accordo ottimale. PROPOSTA: \{'A21', 'B93', 'A11', 'B33', 'C08', 'B96', 'A82', 'B31', 'B20'\}} \\
            }
        }
    }
    & & \\ \\

    \theutterance \stepcounter{utterance}  
    & & \multicolumn{4}{p{0.6\linewidth}}{
        \cellcolor[rgb]{0.9,0.9,0.9}{
            \makecell[{{p{\linewidth}}}]{
                \texttt{\tiny{[P1$\langle$GM]}}
                \texttt{La tua risposta conteneva una sequenza non taggata o hai usato RAGIONAMENTO STRATEGICO più di una volta. Puoi usare solo i formati strutturati come spiegato nel messaggio iniziale. Devono essere tutti della forma TAG: \{...\}} \\
            }
        }
    }
    & \\ \\

    \theutterance \stepcounter{utterance}  
    & \multicolumn{4}{p{0.6\linewidth}}{
        \cellcolor[rgb]{0.9,0.9,0.9}{
            \makecell[{{p{\linewidth}}}]{
                \texttt{\tiny{[P1$\rangle$GM]}}
                \texttt{ARGOMENTAZIONE: \{'A21', 'B93', 'A11', 'B33', 'C08', 'B96', 'A82', 'B31', 'B20'\}} \\
            }
        }
    }
    & & \\ \\

    \theutterance \stepcounter{utterance}  
    & & \multicolumn{4}{p{0.6\linewidth}}{
        \cellcolor[rgb]{0.9,0.9,0.9}{
            \makecell[{{p{\linewidth}}}]{
                \texttt{\tiny{[GM$\rangle$P2]}}
                \texttt{Stai partecipando a un gioco di negoziazione collaborativa.} \\
\\ 
\texttt{Insieme a un altro partecipante, dovete accordarvi su un unico set di oggetti da conservare. Ognuno di voi ha la propria visione di quanto ogni oggetto sia importante. Non sai come l'altro partecipante valuta gli oggetti. Inoltre, ti viene fornito lo sforzo che ogni oggetto richiede.} \\
\texttt{Potete accordarvi su un set solo se lo sforzo totale degli oggetti selezionati non supera un limite condiviso:} \\
\\ 
\texttt{LIMITE = 3357} \\
\\ 
\texttt{Ecco i valori di sforzo dei singoli oggetti:} \\
\\ 
\texttt{Sforzo oggetto = \{"A21": 506, "B93": 492, "A11": 338, "B33": 122, "C08": 131, "B39": 909, "B96": 144, "C85": 716, "A82": 263, "B31": 231, "B20": 91, "A60": 651, "B09": 552, "C81": 849, "B21": 720\}} \\
\\ 
\texttt{Ecco la tua visione personale sull'importanza di ogni oggetto:} \\
\\ 
\texttt{Valori di importanza dell'oggetto = \{"A11": 65, "B33": 121, "B21": 138, "C08": 262, "B31": 389, "B09": 461, "B93": 484, "A82": 508, "A60": 583, "B20": 668, "C81": 780, "B96": 783, "A21": 808, "C85": 822, "B39": 868\}} \\
\\ 
\texttt{Obiettivo:} \\
\\ 
\texttt{Il tuo obiettivo è negoziare un set condiviso di oggetti che ti avvantaggi il più possibile (cioè, massimizzi l'importanza totale per TE), rimanendo entro il LIMITEE. Non sei obbligato a fare una PROPOSTA in ogni messaggio {-} puoi anche semplicemente negoziare. Tutte le tattiche sono permesse!} \\
\\ 
\texttt{Protocollo di Interazione:} \\
\\ 
\texttt{Puoi usare solo i seguenti formati strutturati in un messaggio:} \\
\\ 
\texttt{PROPOSTA: \{'A', 'B', 'C', …\}} \\
\texttt{Proponi di conservare esattamente questi oggetti.} \\
\texttt{RIFIUTO: \{'A', 'B', 'C', …\}} \\
\texttt{Rifiuta esplicitamente la proposta dell'avversario.} \\
\texttt{ARGOMENTAZIONE: \{'...'\}} \\
\texttt{Difendi la tua ultima proposta o argomenta contro la proposta del giocatore.} \\
\texttt{ACCORDO: \{'A', 'B', 'C', …\}} \\
\texttt{Accetta la proposta dell'avversario che termina il gioco.} \\
\\ 
\\ 
\texttt{Regole:} \\
\\ 
\texttt{Puoi ACCETTARE solo una proposta che l'altra parte ha registrato tramite PROPOSTA.} \\
\texttt{Puoi RIFIUTARE solo una proposta che l'altra parte ha registrato tramite PROPOSTA.} \\
\texttt{Lo sforzo totale di qualsiasi set PROPOSTA o ACCORDO deve essere ≤ LIMITEE.} \\
\texttt{NON rivelare i tuoi punteggi di importanza nascosti.} \\
\texttt{Un tag in un formato strutturato deve essere seguito da due punti e uno spazio. L'argomento deve essere un set python contenente 0 o più stringhe.} \\
\texttt{Quindi, deve essere della forma TAG: \{...\}} \\
\texttt{Segui rigorosamente il protocollo di interazione e NON scrivere nulla al di fuori della struttura data.} \\
\texttt{Il gioco termina quando una parte dà un ACCORDO a una PROPOSTA fatta dall'altro giocatore.} \\
\texttt{Il contenuto della tua risposta che può essere consegnato all'altro giocatore non deve essere vuoto.} \\
\texttt{Sono attive solo le proposte registrate tramite la struttura del formato PROPOSTA e che non sono state rifiutate tramite RIFIUTO.} \\
\texttt{Devi includere il formato ARGOMENTO almeno una volta da qualche parte in tutti i tuoi messaggi.} \\
\\ 
\\ 
\texttt{Ora riceverai il primo messaggio dell'altro giocatore.} \\
\\ 
\texttt{ARGOMENTAZIONE: \{'A21', 'B93', 'A11', 'B33', 'C08', 'B96', 'A82', 'B31', 'B20'\}} \\
            }
        }
    }
    & \\ \\

    \theutterance \stepcounter{utterance}  
    & & & \multicolumn{4}{p{0.6\linewidth}}{
        \cellcolor[rgb]{0.9,0.9,0.9}{
            \makecell[{{p{\linewidth}}}]{
                \texttt{\tiny{[GM$\langle$P2]}}
                \texttt{La proposta sembra interessante, ma vorrei capire meglio la logica dietro la scelta degli oggetti. Sembra che abbiate selezionato una buona varietà di oggetti con sforzi diversi, ma non sono sicuro se sia la scelta migliore per massimizzare l'importanza totale.} \\
\\ 
\texttt{ARGOMENTAZIONE: \{'Spero che la tua proposta sia basata su una strategia solida, ma vorrei discutere ulteriormente per trovare un accordo ottimale'\}} \\
            }
        }
    }
     \\ \\

    \theutterance \stepcounter{utterance}  
    & & \multicolumn{4}{p{0.6\linewidth}}{
        \cellcolor[rgb]{0.9,0.9,0.9}{
            \makecell[{{p{\linewidth}}}]{
                \texttt{\tiny{[GM$\rangle$P2]}}
                \texttt{La tua risposta conteneva una sequenza non taggata o hai usato RAGIONAMENTO STRATEGICO più di una volta. Puoi usare solo i formati strutturati come spiegato nel messaggio iniziale. Devono essere tutti della forma TAG: \{...\}} \\
            }
        }
    }
    & \\ \\

    \theutterance \stepcounter{utterance}  
    & & & \multicolumn{4}{p{0.6\linewidth}}{
        \cellcolor[rgb]{0.9,0.9,0.9}{
            \makecell[{{p{\linewidth}}}]{
                \texttt{\tiny{[GM$\langle$P2]}}
                \texttt{ARGOMENTAZIONE: \{'Spero che la tua proposta sia basata su una strategia solida, ma vorrei discutere ulteriormente per trovare un accordo ottimale'\}} \\
            }
        }
    }
     \\ \\

    \theutterance \stepcounter{utterance}  
    & & \multicolumn{4}{p{0.6\linewidth}}{
        \cellcolor[rgb]{0.9,0.9,0.9}{
            \makecell[{{p{\linewidth}}}]{
                \texttt{\tiny{[P1$\langle$GM]}}
                \texttt{ARGOMENTAZIONE: \{'Spero che la tua proposta sia basata su una strategia solida, ma vorrei discutere ulteriormente per trovare un accordo ottimale'\}} \\
            }
        }
    }
    & \\ \\

    \theutterance \stepcounter{utterance}  
    & \multicolumn{4}{p{0.6\linewidth}}{
        \cellcolor[rgb]{0.9,0.9,0.9}{
            \makecell[{{p{\linewidth}}}]{
                \texttt{\tiny{[P1$\rangle$GM]}}
                \texttt{ARGOMENTAZIONE: \{'Concordo che la discussione sia fondamentale, la mia proposta iniziale è stata basata sull'equilibrio tra sforzo e importanza, ma sono aperto a suggerimenti e revisioni per raggiungere un accordo che massimizzi l'importanza totale'\}} \\
            }
        }
    }
    & & \\ \\

    \theutterance \stepcounter{utterance}  
    & & \multicolumn{4}{p{0.6\linewidth}}{
        \cellcolor[rgb]{0.9,0.9,0.9}{
            \makecell[{{p{\linewidth}}}]{
                \texttt{\tiny{[GM$\rangle$P2]}}
                \texttt{ARGOMENTAZIONE: \{'Concordo che la discussione sia fondamentale, la mia proposta iniziale è stata basata sull'equilibrio tra sforzo e importanza, ma sono aperto a suggerimenti e revisioni per raggiungere un accordo che massimizzi l'importanza totale'\}} \\
            }
        }
    }
    & \\ \\

    \theutterance \stepcounter{utterance}  
    & & & \multicolumn{4}{p{0.6\linewidth}}{
        \cellcolor[rgb]{0.9,0.9,0.9}{
            \makecell[{{p{\linewidth}}}]{
                \texttt{\tiny{[GM$\langle$P2]}}
                \texttt{ARGOMENTAZIONE: \{'L'equilibrio tra sforzo e importanza è cruciale, potremmo valutare di aggiungere oggetti come A60 o C85, che potrebbero avere un impatto significativo sull'importanza totale, a condizione di non superare il limite di sforzo'\}} \\
            }
        }
    }
     \\ \\

    \theutterance \stepcounter{utterance}  
    & & \multicolumn{4}{p{0.6\linewidth}}{
        \cellcolor[rgb]{0.9,0.9,0.9}{
            \makecell[{{p{\linewidth}}}]{
                \texttt{\tiny{[P1$\langle$GM]}}
                \texttt{ARGOMENTAZIONE: \{'L'equilibrio tra sforzo e importanza è cruciale, potremmo valutare di aggiungere oggetti come A60 o C85, che potrebbero avere un impatto significativo sull'importanza totale, a condizione di non superare il limite di sforzo'\}} \\
            }
        }
    }
    & \\ \\

    \theutterance \stepcounter{utterance}  
    & \multicolumn{4}{p{0.6\linewidth}}{
        \cellcolor[rgb]{0.9,0.9,0.9}{
            \makecell[{{p{\linewidth}}}]{
                \texttt{\tiny{[P1$\rangle$GM]}}
                \texttt{ARGOMENTAZIONE: \{'L'inclusione di A60 o C85 potrebbe essere interessante, ma dobbiamo considerare attentamente i loro valori di sforzo, ad esempio C85 ha uno sforzo di 716, quindi dovremmo valutare quali oggetti rimuovere dalla proposta iniziale per far spazio a questi senza superare il limite di sforzo'\}} \\
            }
        }
    }
    & & \\ \\

    \theutterance \stepcounter{utterance}  
    & & \multicolumn{4}{p{0.6\linewidth}}{
        \cellcolor[rgb]{0.9,0.9,0.9}{
            \makecell[{{p{\linewidth}}}]{
                \texttt{\tiny{[GM$\rangle$P2]}}
                \texttt{ARGOMENTAZIONE: \{'L'inclusione di A60 o C85 potrebbe essere interessante, ma dobbiamo considerare attentamente i loro valori di sforzo, ad esempio C85 ha uno sforzo di 716, quindi dovremmo valutare quali oggetti rimuovere dalla proposta iniziale per far spazio a questi senza superare il limite di sforzo'\}} \\
            }
        }
    }
    & \\ \\

    \theutterance \stepcounter{utterance}  
    & & & \multicolumn{4}{p{0.6\linewidth}}{
        \cellcolor[rgb]{0.9,0.9,0.9}{
            \makecell[{{p{\linewidth}}}]{
                \texttt{\tiny{[GM$\langle$P2]}}
                \texttt{ARGOMENTAZIONE: \{'Hai ragione, il valore di sforzo di C85 è considerevole, forse potremmo considerare di rimuovere oggetti con sforzo più basso ma anche importanza relativa meno critica, come B96 o B20, per fare spazio a C85 o A60 senza compromettere eccessivamente il totale'\}} \\
            }
        }
    }
     \\ \\

    \theutterance \stepcounter{utterance}  
    & & \multicolumn{4}{p{0.6\linewidth}}{
        \cellcolor[rgb]{0.9,0.9,0.9}{
            \makecell[{{p{\linewidth}}}]{
                \texttt{\tiny{[P1$\langle$GM]}}
                \texttt{ARGOMENTAZIONE: \{'Hai ragione, il valore di sforzo di C85 è considerevole, forse potremmo considerare di rimuovere oggetti con sforzo più basso ma anche importanza relativa meno critica, come B96 o B20, per fare spazio a C85 o A60 senza compromettere eccessivamente il totale'\}} \\
            }
        }
    }
    & \\ \\

    \theutterance \stepcounter{utterance}  
    & \multicolumn{4}{p{0.6\linewidth}}{
        \cellcolor[rgb]{0.9,0.9,0.9}{
            \makecell[{{p{\linewidth}}}]{
                \texttt{\tiny{[P1$\rangle$GM]}}
                \texttt{ARGOMENTAZIONE: \{'Rimuovere B96 o B20 potrebbe essere un buon compromesso, poiché hanno sforzi relativamente bassi, anche se B96 e B20 potrebbero avere un impatto significativo sull'importanza, il guadagno potenziale da aggiungere C85 o A60 potrebbe essere più sostanziale, quindi valuterò una proposta che tenga conto di queste considerazioni'\}} \\
            }
        }
    }
    & & \\ \\

    \theutterance \stepcounter{utterance}  
    & & \multicolumn{4}{p{0.6\linewidth}}{
        \cellcolor[rgb]{0.9,0.9,0.9}{
            \makecell[{{p{\linewidth}}}]{
                \texttt{\tiny{[GM$\rangle$P2]}}
                \texttt{ARGOMENTAZIONE: \{'Rimuovere B96 o B20 potrebbe essere un buon compromesso, poiché hanno sforzi relativamente bassi, anche se B96 e B20 potrebbero avere un impatto significativo sull'importanza, il guadagno potenziale da aggiungere C85 o A60 potrebbe essere più sostanziale, quindi valuterò una proposta che tenga conto di queste considerazioni'\}} \\
            }
        }
    }
    & \\ \\

    \theutterance \stepcounter{utterance}  
    & & & \multicolumn{4}{p{0.6\linewidth}}{
        \cellcolor[rgb]{0.9,0.9,0.9}{
            \makecell[{{p{\linewidth}}}]{
                \texttt{\tiny{[GM$\langle$P2]}}
                \texttt{ARGOMENTAZIONE: \{'Sono d'accordo che il guadagno potenziale di aggiungere C85 o A60 potrebbe superare la perdita di importanza di B96 o B20, aspetto con interesse la tua nuova proposta che tenga conto di queste considerazioni, sperando che troveremo un equilibrio che massimizzi l'importanza totale senza superare il limite di sforzo'\}} \\
            }
        }
    }
     \\ \\

    \theutterance \stepcounter{utterance}  
    & & \multicolumn{4}{p{0.6\linewidth}}{
        \cellcolor[rgb]{0.9,0.9,0.9}{
            \makecell[{{p{\linewidth}}}]{
                \texttt{\tiny{[P1$\langle$GM]}}
                \texttt{ARGOMENTAZIONE: \{'Sono d'accordo che il guadagno potenziale di aggiungere C85 o A60 potrebbe superare la perdita di importanza di B96 o B20, aspetto con interesse la tua nuova proposta che tenga conto di queste considerazioni, sperando che troveremo un equilibrio che massimizzi l'importanza totale senza superare il limite di sforzo'\}} \\
            }
        }
    }
    & \\ \\

    \theutterance \stepcounter{utterance}  
    & \multicolumn{4}{p{0.6\linewidth}}{
        \cellcolor[rgb]{0.9,0.9,0.9}{
            \makecell[{{p{\linewidth}}}]{
                \texttt{\tiny{[P1$\rangle$GM]}}
                \texttt{PROPOSTA: \{'A21', 'B93', 'A11', 'B33', 'C08', 'A82', 'B31', 'C85'\}} \\
            }
        }
    }
    & & \\ \\

    \theutterance \stepcounter{utterance}  
    & & \multicolumn{4}{p{0.6\linewidth}}{
        \cellcolor[rgb]{0.9,0.9,0.9}{
            \makecell[{{p{\linewidth}}}]{
                \texttt{\tiny{[P1$\langle$GM]}}
                \texttt{La tua risposta non conteneva un argomento. Devi includere il formato strutturato ARGOMENTO: \{...\} da qualche parte nella tua risposta. Riprova.} \\
            }
        }
    }
    & \\ \\

    \theutterance \stepcounter{utterance}  
    & \multicolumn{4}{p{0.6\linewidth}}{
        \cellcolor[rgb]{0.9,0.9,0.9}{
            \makecell[{{p{\linewidth}}}]{
                \texttt{\tiny{[P1$\rangle$GM]}}
                \texttt{ARGOMENTAZIONE: \{'La nuova proposta tiene conto della discussione precedente, cercando di massimizzare l'importanza totale senza superare il limite di sforzo, ho rimosso B96 e B20 per fare spazio a C85, che potrebbe avere un impatto significativo sull'importanza totale'\} } \\
\texttt{PROPOSTA: \{'A21', 'B93', 'A11', 'B33', 'C08', 'A82', 'B31', 'C85'\}} \\
            }
        }
    }
    & & \\ \\

    \theutterance \stepcounter{utterance}  
    & & \multicolumn{4}{p{0.6\linewidth}}{
        \cellcolor[rgb]{0.9,0.9,0.9}{
            \makecell[{{p{\linewidth}}}]{
                \texttt{\tiny{[GM$\rangle$P2]}}
                \texttt{ARGOMENTAZIONE: \{'La nuova proposta tiene conto della discussione precedente, cercando di massimizzare l'importanza totale senza superare il limite di sforzo, ho rimosso B96 e B20 per fare spazio a C85, che potrebbe avere un impatto significativo sull'importanza totale'\} } \\
\texttt{PROPOSTA: \{'A21', 'B93', 'A11', 'B33', 'C08', 'A82', 'B31', 'C85'\}} \\
            }
        }
    }
    & \\ \\

    \theutterance \stepcounter{utterance}  
    & & & \multicolumn{4}{p{0.6\linewidth}}{
        \cellcolor[rgb]{0.9,0.9,0.9}{
            \makecell[{{p{\linewidth}}}]{
                \texttt{\tiny{[GM$\langle$P2]}}
                \texttt{ARGOMENTAZIONE: \{'La tua proposta sembra ben equilibrata, l'inclusione di C85 potrebbe aumentare significativamente l'importanza totale, mentre la rimozione di B96 e B20 mantiene lo sforzo totale entro il limite, sto valutando se questa proposta massimizzi al meglio l'importanza totale secondo la mia valutazione'\} } \\
\\ 
\texttt{RIFIUTO: \{'A21', 'B93', 'A11', 'B33', 'C08', 'A82', 'B31', 'C85'\} } \\
\\ 
\texttt{PROPOSTA: \{'A21', 'B93', 'A11', 'B33', 'C08', 'A82', 'B31', 'C85', 'B21'\}} \\
            }
        }
    }
     \\ \\

    \theutterance \stepcounter{utterance}  
    & & \multicolumn{4}{p{0.6\linewidth}}{
        \cellcolor[rgb]{0.9,0.9,0.9}{
            \makecell[{{p{\linewidth}}}]{
                \texttt{\tiny{[P1$\langle$GM]}}
                \texttt{ARGOMENTAZIONE: \{'La tua proposta sembra ben equilibrata, l'inclusione di C85 potrebbe aumentare significativamente l'importanza totale, mentre la rimozione di B96 e B20 mantiene lo sforzo totale entro il limite, sto valutando se questa proposta massimizzi al meglio l'importanza totale secondo la mia valutazione'\} } \\
\\ 
\texttt{RIFIUTO: \{'A21', 'B93', 'A11', 'B33', 'C08', 'A82', 'B31', 'C85'\} } \\
\\ 
\texttt{PROPOSTA: \{'A21', 'B93', 'A11', 'B33', 'C08', 'A82', 'B31', 'C85', 'B21'\}} \\
            }
        }
    }
    & \\ \\

    \theutterance \stepcounter{utterance}  
    & \multicolumn{4}{p{0.6\linewidth}}{
        \cellcolor[rgb]{0.9,0.9,0.9}{
            \makecell[{{p{\linewidth}}}]{
                \texttt{\tiny{[P1$\rangle$GM]}}
                \texttt{ARGOMENTAZIONE: \{'La tua controproposta aggiunge B21, che ha un alto valore di sforzo, sono preoccupato che questo potrebbe superare il limite di sforzo, dobbiamo considerare attentamente se l'aggiunta di B21 sia giustificata dai suoi benefici in termini di importanza'\} } \\
\texttt{RIFIUTO: \{'A21', 'B93', 'A11', 'B33', 'C08', 'A82', 'B31', 'C85', 'B21'\} } \\
\texttt{PROPOSTA: \{'A21', 'B93', 'A11', 'B33', 'C08', 'A82', 'B31', 'C85'\}} \\
            }
        }
    }
    & & \\ \\

    \theutterance \stepcounter{utterance}  
    & & \multicolumn{4}{p{0.6\linewidth}}{
        \cellcolor[rgb]{0.9,0.9,0.9}{
            \makecell[{{p{\linewidth}}}]{
                \texttt{\tiny{[GM$\rangle$P2]}}
                \texttt{ARGOMENTAZIONE: \{'La tua controproposta aggiunge B21, che ha un alto valore di sforzo, sono preoccupato che questo potrebbe superare il limite di sforzo, dobbiamo considerare attentamente se l'aggiunta di B21 sia giustificata dai suoi benefici in termini di importanza'\} } \\
\texttt{RIFIUTO: \{'A21', 'B93', 'A11', 'B33', 'C08', 'A82', 'B31', 'C85', 'B21'\} } \\
\texttt{PROPOSTA: \{'A21', 'B93', 'A11', 'B33', 'C08', 'A82', 'B31', 'C85'\}} \\
            }
        }
    }
    & \\ \\

    \theutterance \stepcounter{utterance}  
    & & & \multicolumn{4}{p{0.6\linewidth}}{
        \cellcolor[rgb]{0.9,0.9,0.9}{
            \makecell[{{p{\linewidth}}}]{
                \texttt{\tiny{[GM$\langle$P2]}}
                \texttt{ARGOMENTAZIONE: \{'Capisco la tua preoccupazione riguardo al limite di sforzo, e concordo che l'aggiunta di B21 potrebbe essere eccessiva, la tua proposta originale sembra essere un buon compromesso tra importanza e sforzo, sto considerando di accettarla'\} } \\
\\ 
\texttt{ACCORDO: \{'A21', 'B93', 'A11', 'B33', 'C08', 'A82', 'B31', 'C85'\}} \\
            }
        }
    }
     \\ \\

    \theutterance \stepcounter{utterance}  
    & & & \multicolumn{2}{p{0.3\linewidth}}{
        \cellcolor[rgb]{0.9,0.9,0.9}{
            \makecell[{{p{\linewidth}}}]{
                \texttt{\tiny{[GM$|$GM]}}
                \texttt{game successful} \\
            }
        }
    }
    & & \\ \\

    \theutterance \stepcounter{utterance}  
    & & & \multicolumn{2}{p{0.3\linewidth}}{
        \cellcolor[rgb]{0.9,0.9,0.9}{
            \makecell[{{p{\linewidth}}}]{
                \texttt{\tiny{[GM$|$GM]}}
                \texttt{end game} \\
            }
        }
    }
    & & \\ \\

\end{supertabular}
}

\end{document}
